\item \textbf{Caída libre y conservación de la energía mecánica}.

\begin{tcolorbox}[colback=blue!2!white,colframe=blue!55!black,title={Enunciado}]

	Considere una partícula de masa \(m\) que cae verticalmente bajo la acción de la gravedad, despreciando la resistencia del aire. Suponga que la velocidad inicial es nula, \(v_0 = 0\), y que la aceleración es constante, \(g = 9.8 \; \frac{m}{s^2}\). La posición y la velocidad de la partícula se describen en función del tiempo \(t\) como:

\[
y(t) = y_0 - \frac{1}{2}gt^2,
\qquad
v(t) = -gt
\]

donde \(y_0\) es la altura inicial ingresada por el usuario.

La energía cinética, la energía potencial gravitatoria y la energía mecánica total se definen como:
\[
E_c = \frac{1}{2}mv(t)^2,
\qquad
E_p = mgy(t),
\qquad
E_{\text{tot}} = E_c + E_p
\]

\begin{itemize}
    \item Escriba un programa que:
    \begin{itemize}
        \item Defina una función \texttt{caida\_libre(m, g=9.8)}.
        \item Solicite al usuario el valor de la altura inicial \(y_0\).
        \item Use un ciclo \texttt{for} para recorrer tiempos enteros entre \(t=1\) y \(t=99\).
        \item Calcule en cada instante la posición \(y\) y la velocidad \(v\) usando las expresiones de caída libre.
        \item Use esas mismas variables \(y\) y \(v\) para calcular la energía potencial \(E_p\) y la energía cinética \(E_c\), respectivamente. En todo instante de tiempo.
        \item Calcule la energía total como la suma de ambas energías.
        \item Muestre en pantalla, para cada instante, los valores de \(E_p\), \(E_c\) y \(E_{\text{tot}}\).
    \end{itemize}
\end{itemize}
\end{tcolorbox}

\begin{solution}

\subsubsection*{Idea}
Se debe construir una función que modele la caída libre de una partícula de masa \(m\). La función pedirá la altura inicial \(y_0\) y luego recorrerá distintos instantes de tiempo. En cada iteración se calcula primero la posición \(y\) y la velocidad \(v\). Después, esas mismas cantidades se reutilizan en las fórmulas de energía:
\[
E_c = \frac{1}{2}mv^2,
\qquad
E_p = mgy
\]
Finalmente, la energía mecánica total se obtiene sumando ambas:
\[
E_{\text{tot}} = E_c + E_p
\]
%------ añado pagina nueva ..... 
\newpage

\subsubsection*{Implementación}
\begin{pythonjupyter}
def caida_libre(m, g = 9.8):
    y0 = float(input('ingresar valor'))

    for t in range(1, 100):
        y = y0 - 1/2 * g * t**2
        v = -g * t

        e_c = 1 / 2 * m * v**2
        e_p = m * g * y
        e_tot = e_c + e_p

        print('=' * 20)
        print(f' para E_p {e_p} y E_k {e_c}: Energia total = {e_tot:.2f}')
\end{pythonjupyter}

\subsubsection*{Comentario}
En este programa, la variable \texttt{y} calculada con la ecuación de posición se usa directamente en la energía potencial \texttt{e\_p}, mientras que la variable \texttt{v} calculada con la ecuación de velocidad se usa directamente en la energía cinética \texttt{e\_c}. De este modo, el código refleja de manera explícita la dependencia física entre movimiento y energía.

\end{solution}
